\begin{center}
    {\LARGE \textbf{2014无机化学}} \\[2ex]
    {\large 时量180分钟 \quad 满分100分}
\end{center}

\vspace{0.5cm}
\noindent\textbf{注意事项:}
\begin{enumerate}[label=\arabic*., parsep=0pt, itemsep=0pt]
    \item 答题前,考生先将自己的姓名、学号、考场号、座位号填写在答题卡上,将条形码准确粘贴在考生信息条形码粘贴区。
    \item 考试结束后,将本试卷与答题纸一并收回。
    \item 回答各个问题时,将答案写在答题纸上,写在本试卷上无效。
\end{enumerate}

\vspace{0.5cm}
\noindent\textbf{一、写出下列反应的化学方程式，并叙述实验现象（25 pts.）}

\begin{enumerate}[label=\arabic*., start=1, itemsep=1.5ex]
    \item 向\ce{Hg(NO3)2}溶液中缓慢滴加\ce{KI}溶液至过量。
    \item 使砷化氢通过一玻璃管，并在玻璃管的入口处加热。
    \item 向\ce{NiSO4}溶液中加入过量\ce{NaOH}，充分反应后滴加氯水。
    \item 向\ce{AgNO3}溶液中滴加\ce{Na2S2O3}溶液直至过量，再用盐酸酸化。
    \item 在重铬酸钾\ce{K2Cr2O7}的酸性溶液中加入戊醇，再加入少量\ce{H2O2}后充分振荡。
\end{enumerate}

\vspace{0.5cm}
\noindent\textbf{二、回答下列问题（70 pts.）}

\begin{enumerate}[label=\arabic*., start=1, itemsep=1.5ex]
    \item 给出具有下列电子构型的元素的名称、符号、在周期表中的位置：
    \begin{enumerate}[label=(\arabic*), itemsep=1ex]
        \item $3d^54s^2$
        \item $4d^75s^1$
        \item $5s^25p^1$
        \item $5d^{10}6s^2$
        \item $4s^24p^4$
    \end{enumerate}
    \item 指出下列分子或离子的几何构型和中心的杂化方式。
    \begin{enumerate}[label=(\arabic*), itemsep=1ex]
        \item \ce{H3PO2}
        \item \ce{XeOF4}
        \item \ce{Fe(CO)5}
        \item \ce{[Ni(CN)4]^{2-}}
        \item \ce{BrF3}
    \end{enumerate}
    \item 指出元素\ce{C}，\ce{N}，\ce{O}，\ce{P}和\ce{S}中，第一电离能最大和最小的元素及第一电子亲合能最大和最小的元素，并简述理由。
    \item 判断离子：\ce{H^-}、\ce{H_2^+}、\ce{H_2^-}、\ce{H_2^{2-}}哪些可以存在？哪些不能存在？并说明理由。
    \item 比较\ce{N}的四种氢化物：\ce{HN3}、\ce{NH3}、\ce{NH2OH}和\ce{N2H4}的碱性强弱，并解释原因。
    \item 试画图表示：焦硅酸根\ce{Si2O7^{4-}}，链三聚硅酸根和环六聚硅酸根分子内原子间的键联关系。写出链$n$聚和环$n$聚硅酸根的化学式。
    \item 讨论下列卤化物的熔点变化规律，并说明原因。
    \[
    \begin{array}{lccccc}
     & \ce{TiF4} & \ce{TiCl4} & \ce{TiBr4} & \ce{TiI4} \\
    \text{熔点} & 284^\circ\text{C} & -24^\circ\text{C} & 38^\circ\text{C} & 150^\circ\text{C}
    \end{array}
    \]
    \item 用两种方法鉴别下列各对物质。
    \begin{enumerate}[label=(\arabic*), itemsep=1ex]
        \item \ce{SbCl3}和\ce{BiCl3}
        \item \ce{Na2SO4}和\ce{Na2SO3}
    \end{enumerate}
    \item 比较下列物质的稳定性，简要说明理由。
    \begin{enumerate}[label=(\arabic*), itemsep=1ex]
        \item \ce{Na2CO3}，\ce{MgCO3}，\ce{K2CO3}，\ce{MnCO3}
        \item \ce{HgF4^{2-}}，\ce{HgCl4^{2-}}，\ce{HgBr4^{2-}}，\ce{HgI4^{2-}}
    \end{enumerate}
    \item 分别用价键理论和晶体场理论讨论\ce{[CuCl4]^{2-}}的正方形构型。
    \item 计算八面体强场中$d^7$组态和四面体弱场$d^8$组态的CFSE，用简图表示$d$电子在晶体场中的排布情况。
    \item 已知：
    \[
    \begin{array}{ll}
    \ce{Cu^+ + e^- = Cu} & E^\ominus = 0.52V \\
    \ce{Cu(NH3)_2^+ + e^- = 2NH3 + Cu} & E^\ominus = -0.10V
    \end{array}
    \]
    设计原电池反应，计算298K时\ce{Cu(NH3)_2^+}的稳定常数。
\end{enumerate}

\vspace{0.5cm}
\noindent\textbf{三、制备与分离（25 pts.）}

\begin{enumerate}[label=\arabic*., start=1, itemsep=1.5ex]
    \item 智利硝石中的碘酸钠是重要的碘资源，试叙述以此为原料提取单质碘的方法。
    \item 1986年美国化学家克里斯特（Christe）使用\ce{KMnO4}，\ce{HF}，\ce{KF}，\ce{H2O2}和\ce{SbCl5}为原料，用化学法制得单质\ce{F2}。试用反应方程式表示各步反应过程。
    \item 四瓶标签已经脱落的白色固体试剂：磷酸氢二钠、亚磷酸钠、偏磷酸钠、焦磷酸钠，试通过实验加以鉴别。写出简单步骤、实验现象和反应方程式。
    \item 设计方案分离混合溶液中的离子。
    \begin{enumerate}[label=(\arabic*), itemsep=1ex]
        \item \ce{Ag^+}、\ce{Mn^{2+}}、\ce{Fe^{3+}}、\ce{Ni^{2+}}、\ce{Co^{2+}}
        \item \ce{Sn^{2+}}、\ce{Cu^{2+}}、\ce{Al^{3+}}、\ce{Zn^{2+}}
    \end{enumerate}
\end{enumerate}

\vspace{0.5cm}
\noindent\textbf{四、推断（30 pts.）}

\begin{enumerate}[label=\arabic*., start=1, itemsep=1.5ex]
    \item 化合物\ce{A}与水作用生成白色沉淀\ce{B}。用\ce{NaClO}处理\ce{B}，得到土黄色沉淀\ce{C}。将\ce{C}与少量酸性\ce{MnSO4}溶液混合，\ce{C}溶液，溶液变为红色。用煤气灯加热\ce{A}得到黄色固体\ce{D}和棕色气体\ce{E}。将气体\ce{E}通过\ce{NaOH}溶液后，气体体积仅剩约20\%，将\ce{A}与\ce{SnCl2}混合后加入\ce{NaOH}溶液，生成黑色沉淀\ce{F}。
    试确定字母\ce{A}，\ce{B}，\ce{C}，\ce{D}，\ce{E}，\ce{F}所代表物质的化学式，给出相关各步骤的化学反应方程式。
    \item 有一固体混合物可能含有\ce{FeCl3}，\ce{NaNO2}，\ce{Ca(OH)2}，\ce{AgNO3}，\ce{CuCl2}，\ce{NaF}，\ce{NH4Cl}七种物质中的若干种。若将此混合物加水后，可得白色沉淀和无色溶液，在此无色溶液中加入\ce{KSCN}没有变化，无色溶液可使酸化后\ce{KMnO4}溶液紫色褪去；将无色溶液加热有气体放出。另外，白色沉淀可溶于\ce{NH3\cdot H2O}中。
    试根据上述现象，判断在混合物中，哪些物质肯定存在？哪些肯定不存在？哪些可能存在？简述理由，并用方程式表示。
    \item 绿色化合物\ce{A}溶于水，通入\ce{SO2}至溶液为中性，得到近无色的溶液\ce{B}和黑色沉淀\ce{C}，继续通入\ce{SO2}则沉淀\ce{C}逐渐溶直至全部消失，用稀\ce{H2SO4}处理\ce{A}则得到\ce{D}的紫色溶液和沉淀\ce{C}，再加入过量的\ce{H2O2}则沉淀溶解形成溶液\ce{B}。将\ce{D}的溶液用\ce{HNO3}酸化后加入\ce{Cr(NO3)3}，得到橙色溶液\ce{E}，向\ce{E}中加入\ce{NaOH}，则溶液变为黄色，同时有白色沉淀\ce{F}生成。\ce{D}与冷浓硫酸作用，可得绿色油状化合物\ce{G}，\ce{G}不稳定，易发生爆炸分解生成\ce{C}和无色气体\ce{H}。
    试给出\ce{A}，\ce{B}，\ce{C}，\ce{D}，\ce{E}，\ce{F}，\ce{G}和\ce{H}的化学式，并完成各步的化学反应方程式。
\end{enumerate}